% Generated by scripts/make_tables.py from data/report_data.json.
\begingroup
\footnotesize
\setlength{\tabcolsep}{5pt}
\renewcommand{\arraystretch}{1.2}
\begin{tabularx}{\linewidth}{@{}p{0.06\linewidth}p{0.15\linewidth}p{0.31\linewidth}>{\raggedright\arraybackslash}X@{}}
\toprule
题号 & 主要问题 & 已观察现象 & 核查证据与待验证改进 \\
\midrule
\textbf{F02} & 生成／传参偏移 & 将集成模型测试集 3.57\% 误写为单模型验证集结果。 & P06，第 2、7 页区分两种设置；建议保留原题限定，并逐项绑定模型、数据划分与指标。 \\
\addlinespace[4pt]
\textbf{F07} & Agent 执行失败 & 事实问题误走元数据工具，字符串拼接 TypeError 导致无答案、无引用。 & P12，第 1 页包含所问设置；建议修正意图判定、规范字段类型，并保留真实失败状态。 \\
\addlinespace[4pt]
\textbf{F16} & 检索证据缺失 & 未返回 ViT 数据集规模所在页，最终拒答三组类别数／图像数。 & P03，第 4 页有完整数据；建议约束目标来源与章节，再检验召回和答案完整性。 \\
\addlinespace[4pt]
\textbf{C14} & 来源匹配／生成 & Vim 匹配成另一篇论文；又将 ResNet 的 ImageNet 指标误归 CIFAR-10。 & P02，第 2、5 页及 ResNet 第 7 页可核对；建议使用确定的题名—文件映射，并按数据集组织证据。 \\
\addlinespace[4pt]
\textbf{S10} & 任务分解／状态 & 三篇论文被交给双篇比较工具；返回错误信息，却以 success=True 结束。 & 三篇原文第 1 页均有相应内容；建议逐篇抽取后汇总，并将工具失败状态向上传递。 \\
\addlinespace[4pt]
\textbf{R08} & 推理归因错误 & 以色彩基线缺少端到端评估，解释复合规则 0/102 不能算独立验证误报率。 & P14，第 8、11 页说明同集设计偏差及冻结后独立验证；建议区分评估对象、设计集与验证集。 \\
\bottomrule
\end{tabularx}
\par\vspace{4pt}{\footnotesize 案例由冻结的 AI 辅助评分与实际结果归纳；P 编号对应文献表，页码均为 PDF 物理页。改进为待实施、待复测建议，不代表已取得优化后结果。\par}
\endgroup
